\section{Proofs}
\label{app:proofs}

Proposition~\ref{prop:axioms} (axiom verification) is proved in the main text
(Section~\ref{sec:axioms}). This appendix provides full proofs for the
remaining results.

% ---------------------------------------------------------------------------
% A.1: Tractability
% ---------------------------------------------------------------------------

\subsection{Proof of Proposition~\ref{prop:tractability} (Polynomial-Time Computation)}
\label{app:proof_tractability}

\begin{proof}
We describe the algorithm explicitly, then verify its correctness and
complexity.

\paragraph{Algorithm.}

\emph{Input:} A finite capability poset $(\C, \leq)$ with $|\C| = n$ nodes,
weight function $w$, and a set $G \subseteq \C$.

\emph{Output:} $\volL(G)$.

\begin{enumerate}
\item \textbf{Topological sort.} Compute a topological ordering of $(\C,
  \leq)$ from $\bot$ upward. Cost: $O(n + e)$ where $e$ is the number of
  edges in the Hasse diagram. Since $e \leq n \cdot k$ (each node has at most
  $k$ covers), this is $O(n \cdot k)$.

\item \textbf{Compute M\"{o}bius function and independent weights.} Compute
  $\mu(d', d)$ for all pairs $d' \leq d$ by processing in topological order.
  For each node $d$ and each $d' \leq d$:
  $\mu(d', d) = -\sum_{d' \leq d'' < d} \mu(d', d'')$.
  Then $\iw(d) = \sum_{d' \leq d} \mu(d', d) \cdot w(d')$.
  There are $O(n^2)$ comparable pairs, each requiring a sum over up to
  $O(n)$ intermediate elements: total $O(n^3)$. For tree-structured posets
  (no diamonds), $\mu$ is nonzero only between covers, reducing to
  $O(n \cdot k)$.

\item \textbf{Compute downward closure.} Initialize $\dc G = \emptyset$.
  For each $d \in G$, perform a depth-first traversal downward in the Hasse
  diagram, adding all visited nodes to $\dc G$. Each node is visited at most
  once (mark visited). Total: $O(n)$ (each node in $\C$ is visited at most
  once across all traversals).

\item \textbf{Sum.} Compute $\volL(G) = \sum_{d \in \dc G} \iw(d)$. Cost:
  $O(|\dc G|) \leq O(n)$.
\end{enumerate}

\emph{Total cost:} $O(n^3)$ for general posets, dominated by the full
M\"{o}bius computation in step~2. For tree-structured posets (each node
has at most one cover), $\mu$ is nonzero only between covers, reducing the
cost to $O(n \cdot k)$ where $k$ is the maximum in-degree.

\paragraph{Correctness.} Step~2 computes $\iw(d)$ by M\"{o}bius inversion on
the poset $(\C, \leq)$, which is well-defined for any finite poset
\citep{rota1964foundations}. Step~3 computes $\dc G$ by the standard downward
closure algorithm on a DAG. Step~4 sums non-negative terms (by the design
constraint $\iw(d) \geq 0$) over the closure, yielding $\volL(G)$ by
Definition~\ref{def:poset_measure}.
\end{proof}

% ---------------------------------------------------------------------------
% A.2: Sign-Correctness
% ---------------------------------------------------------------------------

\subsection{Proof of Proposition~\ref{prop:sign_correctness} (Upgraded Sign-Correctness)}
\label{app:proof_sign_correctness}

\begin{proof}
The structure parallels Appendix~D.8 of \citet{lasser2026gfm}, with the key
difference that $R_k(\pi)$ is now real-valued rather than ternary.

\paragraph{Setup.} Decompose the true volume change as
$\Delta \volL(G) = \Delta_I + \Delta_C$, where
$\Delta_I = \sum_k \sum_{d \in D_k(\pi)} \iw(d)$ sums over all agents'
gained and lost capabilities (the individual component) and $\Delta_C$ is the
cooperative component.

The estimator observes:
\[
\Deltahat(\pi) = \sum_{k \in \mathcal{O}} T_k \cdot R_k(\pi)
\]
where $R_k(\pi) = \sum_{d \in D_k^+(\pi)} \iw(d) - \sum_{d \in D_k^-(\pi)}
\iw(d)$ (Equation~\eqref{eq:weighted_report}).

Define the individual estimation error $e_I = \Delta_I - \hat{\Delta}_I$ and
the cooperative estimation error $e_C = \Delta_C - \hat{\Delta}_C$.

\paragraph{Error decomposition.} Under conditions 1--2 (representative sample,
accurate trust model), the individual error satisfies $|e_I| \leq \delta \cdot
|\Delta_I|$ for a sampling error ratio $\delta$ controlled by the
observation coverage $|\mathcal{O}|$ and trust calibration.

Because $R_k$ carries real-valued weights in the same units as $\iw(d)$, the
cooperative error $e_C$ accounts for cooperative capabilities whose
participants are outside $\mathcal{O}$. For cooperative capabilities with at
least one participant in $\mathcal{O}$, the participant's report includes the
cooperative change (Section~2.1 of the main text), so $e_C$ is bounded by the
volume of cooperative capabilities involving only unobserved agents.

\paragraph{Sign-correctness.} We need $\sign(\Deltahat) = \sign(\Delta
\volL)$.

\emph{Same-sign case:} When $\Delta_I$ and $\Delta_C$ have the same sign,
$|\Delta \volL| = |\Delta_I| + |\Delta_C| \geq |\Delta_I|$. The total error
$|e_I + e_C| \leq |e_I| + |e_C| < (\delta + 1 - \epsilon)|\Delta_I| \leq
|\Delta \volL|$ whenever $\delta < \epsilon$. Sign is preserved.

\emph{Opposite-sign case:} When $\Delta_I$ and $\Delta_C$ have opposite
signs, $|\Delta \volL| = \bigl||\Delta_I| - |\Delta_C|\bigr|$. Sign
preservation requires $|e_I + e_C| < \bigl||\Delta_I| - |\Delta_C|\bigr|$.
This holds when $|\Delta_C| < (\epsilon - \delta)|\Delta_I|$.

\paragraph{Improvement over ternary estimator.} With ternary signals, an
agent that gains a high-weight capability and loses a low-weight one reports
$R_k = 0$ (one gain, one loss), losing the magnitude asymmetry. With
real-valued signals, the same agent reports
$R_k = \iw(d_{\mathrm{gained}}) - \iw(d_{\mathrm{lost}})$, preserving the
net direction. This strictly reduces $|e_I|$ for any action that produces
asymmetric capability changes, widening the sign-correctness margin.
\end{proof}

% ---------------------------------------------------------------------------
% A.3: Leverage Amplifies Cooperative Scoring
% ---------------------------------------------------------------------------

\subsection{Proof of Proposition~\ref{prop:leverage_cooperation} (Leverage Amplifies Cooperation)}
\label{app:proof_leverage_cooperation}

\begin{proof}
Let $d_c \in G^{\mathrm{coop}}_{ij}(t+1) \setminus G^{\mathrm{coop}}_{ij}(t)$
be a new cooperative capability created by joint action $\pi_{ij}$, entering
the poset as an atom. The leverage-adjusted score
(Equation~\eqref{eq:leverage_score}) for the cooperative action is:
\begin{align*}
S(\pi_{ij}) &= \Deltahat(\pi_{ij}) + \alpha \sum_{d_c \in D^+(\pi_{ij})}
  \hat{\lambda}(d_c) \cdot \iw(d_c)
\end{align*}

Since $d_c$ requires both $a_i$ and $a_j$, every time step involving $d_c$
also involves both agents' meta-capabilities. The leverage estimator attributes
the excess $\Deltahat$ in those time steps to $d_c$: since either partner's
meta-capability alone produces at least $\lambda_{\min}$ excess growth,
$\hat{\lambda}(d_c) \geq \lambda_{\min}$.

For new atoms, $\iw(d_c) = w(d_c)$ (no M\"{o}bius subtraction), so the
leverage bonus is at least $\alpha \lambda_{\min} \cdot \iw(d_c) = \alpha
\lambda_{\min} \cdot \Delta \volL(G^{\mathrm{coop}}_{ij})$ (summing over all
new atoms). Therefore:
\[
S(\pi_{ij}) \geq \Deltahat(\pi_{ij}) + \alpha \lambda_{\min} \cdot
\Delta \volL(G^{\mathrm{coop}}_{ij})
\]

Because leverage operates as a ranking signal rather than a measure
modification, this result does not require verifying non-negativity of any
leverage-weighted independent weights. The base measure $\volL$ is unchanged;
only the action-selection criterion is modified.
\end{proof}

% ---------------------------------------------------------------------------
% A.4: Bootstrap Convergence
% ---------------------------------------------------------------------------

\subsection{Proof of Proposition~\ref{prop:bootstrap_convergence} (Poset Growth)}
\label{app:proof_bootstrap}

\begin{proof}
We verify each part of the proposition.

\paragraph{(a) Node set monotonicity.} The poset operations (add node, add
edge, update benchmark) never remove nodes or edges. Therefore $\C(t)
\subseteq \C(t+1)$ for all $t$.

\paragraph{(b) Capability addition is $\volL$-non-decreasing.} Adding a new
capability $d$ to $G$ without changing the subsumption structure gives
$G(t+1) \supseteq G(t)$ and $\dc G(t+1) \supseteq \dc G(t)$. Since all
independent weights are unchanged and non-negative, $\volL(G(t+1)) \geq
\volL(G(t))$ by M3.

\paragraph{(c) Edge discovery can decrease $\volL$.} Suppose $d_1, d_2 \in G$
were initially incomparable, with $\volL(\{d_1, d_2\}) = w(d_1) + w(d_2)$ by
M4 (poset-disjointness). Discovering $d_1 \leq d_2$ triggers M\"{o}bius
recomputation: $\iw(d_2)$ decreases from $w(d_2)$ to $w(d_2) - w(d_1)$, and
$\volL(\{d_1, d_2\}) = w(d_1) + (w(d_2) - w(d_1)) = w(d_2) < w(d_1) +
w(d_2)$. The decrease reflects the correction of an overcount: the actor
previously credited $d_1$ and $d_2$ independently when $d_1$'s content was
already captured by $d_2$.

\paragraph{(d) Benchmark refinement: position-dependent.} Increasing
$\smax(d)$ increases $w(d)$. The effect on $\volL(G)$ is determined by the
coefficient of $w(d)$ in the expanded double sum. Write $\volL(G) =
\sum_{d'' \in \dc G} \sum_{d' \leq d''} \mu(d', d'') \cdot w(d')$ and
rearrange by $d'$: the coefficient of $w(d)$ is
$c(d) = \sum_{d'' \geq d,\, d'' \in \dc G} \mu(d, d'')$.

\emph{Case 1: $d$ is maximal in $\dc G$.} Then $c(d) = \mu(d, d) = 1$, so
$\volL$ increases by $\Delta w(d)$.

\emph{Case 2: $d$ has a single ancestor $d'$ in $\dc G$.} Then $c(d) =
\mu(d, d) + \mu(d, d') = 1 + (-1) = 0$, so $\volL$ is unchanged.

\emph{Case 3: $d$ is a shared prerequisite of $m > 1$ elements in $\dc G$.}
Counterexample: let $d < y$ and $d < z$ with $G = \{y, z\}$. Then
$\dc G = \{d, y, z\}$ and $c(d) = \mu(d, d) + \mu(d, y) + \mu(d, z) =
1 - 1 - 1 = -1 < 0$. Increasing $w(d)$ \emph{decreases} $\volL(G)$. This is
epistemically correct: the shared prerequisite's information content is
counted once in $\volL$ regardless of how many capabilities depend on it, and
increasing the prerequisite's weight increases the subtraction in multiple
ancestors' independent weights, reducing total volume.

In general, $c(d)$ can take any integer value depending on the structure of
$\dc G$ above $d$. Benchmark refinement is therefore not monotone on
branched posets.
\end{proof}
